\chapter{The Big Idea: Tensor-Field Simulation}
\label{ch:bigidea}

\chapterorient{We have seen what the simulator computes (\cref{ch:targets}) and
the stages every run passes through (\cref{ch:lifecycle}). This chapter is about
the single idea that organizes the whole book---the lens through which we will
read every operator, every solver, every driver. It is a small idea, and a
familiar one once you see it: \emph{a simulation is a sequence of transformations
of immutable arrays of numbers.} Not a program that scribbles over a block of
memory, but a chain of pure functions, each taking a field and returning a new
one. By the end you will understand why we choose this lens, what it buys us, and
the three questions it teaches us to ask of any computation.}

\section{A field is a tensor}

Every analysis in \cref{ch:targets} solves for a \emph{field}: a quantity that
takes a value at every point of space. The electrostatic analysis solves for the
potential $V$, one number at each point. The magnetostatic analysis solves for
the vector potential $\vA$, three numbers at each point. The eigenmode analysis
solves for a mode field $\vE$, again a vector at each point.

A computer cannot store a value at \emph{every} point---there are infinitely
many. So the first thing any simulator does is \emph{discretize}: it replaces the
continuous field with a finite list of numbers, the \emph{degrees of freedom}
(DoFs), from which the field everywhere can be reconstructed. How that list is
chosen---which points, which basis functions---is the business of the finite
element method, and we build it carefully in \cref{part:fieldtheory}. For now,
only the shape of the answer matters: \emph{a discretized field is a finite,
ordered collection of numbers}. If a mesh has $N$ degrees of freedom, the
potential becomes a list of $N$ real numbers; the $\vE$-field becomes a list of
$N$ numbers too (one per edge degree of freedom).

Such an ordered, rectangular collection of numbers is exactly what we call a
\emph{tensor}. A list of $N$ numbers is a tensor of shape $[N]$. A grid of $H$
rows and $W$ columns is a tensor of shape $[H, W]$. A field that stores a
3-vector at each cell of an $H \times W$ grid is a tensor of shape $[H, W, 3]$.
The numbers $N$, or $H, W, 3$, are the \emph{axes} of the tensor; each axis has a
length, and we will increasingly give each axis a \emph{name} so we never lose
track of what it counts.

\begin{keyidea}
A discretized field \emph{is} a tensor---a finite, named-axis array of numbers.
The whole simulator, from this height, is a machine that consumes tensors
(the mesh, the materials, the excitation) and produces tensors (the solution
field) and finally small tensors (the capacitance matrix, the table of
frequencies). To understand the simulator is to understand the transformations
that carry one tensor to the next.
\end{keyidea}

\begin{figure}[t]
\centering
\begin{tikzpicture}[scale=1.0]
  % --- left: a field over a mesh, as Tensor[N] ---
  \begin{scope}[shift={(-3.4,0)}]
    % little triangulated patch
    \draw[simGray, semithick] (-0.9,-0.7) -- (0.9,-0.7) -- (0.6,0.7) -- (-0.7,0.6) -- cycle;
    \draw[simGray, semithick] (-0.9,-0.7) -- (0.6,0.7);
    \draw[simGray, semithick] (0.9,-0.7) -- (-0.7,0.6);
    \foreach \p in {(-0.9,-0.7),(0.9,-0.7),(0.6,0.7),(-0.7,0.6),(-0.05,-0.0)}{
      \fill[simRust] \p circle (1.4pt);}
    \node[font=\scriptsize\itshape, simGray] at (0,-1.15) {a field on a mesh};
  \end{scope}
  \node[font=\Large] at (-1.7,0) {$\rightsquigarrow$};
  % the Tensor[N] box
  \node[opbox, minimum width=30mm, minimum height=9mm] (tN) at (0.1,0)
    {\lstinline[style=l4style]|Tensor[N]|};
  \node[font=\scriptsize\itshape, simGray, below=1mm of tN]
    {$N$ degrees of freedom, one axis};
  % --- right: a multi-axis tensor Tensor[H, W, 3] ---
  \begin{scope}[shift={(4.6,0)}]
    \def\dx{0.16}\def\dy{0.10}
    % a small 3-deep stack of HxW grids to suggest axes
    \foreach \k/\c in {0/simBgGray,1/simBgBlue,2/simBgRust}{
      \begin{scope}[shift={(\k*\dx,\k*\dy)}]
        \fill[\c, draw=simGray] (-0.8,-0.55) rectangle (0.8,0.55);
        \draw[simGray!60] (-0.8,-0.18) -- (0.8,-0.18);
        \draw[simGray!60] (-0.8,0.18) -- (0.8,0.18);
        \draw[simGray!60] (-0.27,-0.55) -- (-0.27,0.55);
        \draw[simGray!60] (0.27,-0.55) -- (0.27,0.55);
      \end{scope}}
    \node[font=\scriptsize, simInk] at (-1.15,0.0) {$H$};
    \node[font=\scriptsize, simInk] at (0.0,-0.85) {$W$};
    \node[font=\scriptsize, simRust] at (1.15,0.75) {$3$};
    \node[font=\scriptsize\ttfamily, simInk, below=4mm of {(0,-0.55)}] (lab)
      at (0.1,-1.0) {Tensor[H, W, 3]};
  \end{scope}
\end{tikzpicture}
\caption[A tensor as a labeled box of axes]{Two tensors, drawn as labeled boxes
of axes. \emph{Left:} a scalar field on a mesh of $N$ degrees of freedom is a
one-axis tensor, written \lstinline[style=l4style]|Tensor[N]|. \emph{Right:} a
field storing a 3-vector at every cell of an $H\times W$ grid is a three-axis
tensor \lstinline[style=l4style]|Tensor[H, W, 3]|, with named axes $H$, $W$, and
the length-3 component axis. Reading and reasoning about these shapes is a skill
we develop in full in \cref{ch:tensors}; here we need only the picture.}
\label{fig:tensorbox}
\end{figure}

\Cref{fig:tensorbox} shows the two pictures we will lean on. On the left, a
scalar field over a mesh, collapsed into a single-axis tensor
\lstinline[style=l4style]|Tensor[N]|; on the right, a small multi-axis tensor
\lstinline[style=l4style]|Tensor[H, W, 3]| whose three named axes remind us that
a tensor is not just ``a big array'' but an array \emph{with structure we can
name}. We will treat shapes seriously---they are a kind of type, and they catch
mistakes the way ordinary types do---beginning in \cref{ch:tensors}. In this
chapter we keep them at arm's length and let the picture do the work.

\section{Two ways to evolve a state}

Now suppose we have our field as a tensor and we want to advance it: collide and
stream a fluid, take a time step, apply one iteration of a solver. There are two
very different mental models for ``advancing,'' and the choice between them is the
true subject of this chapter.

\subsection{The picture you probably have: in-place mutation}

If you have written numerical code in C or Fortran---or read any classical
simulation textbook---you carry a particular picture of how a field evolves. You
allocate an array. You loop over its entries. You \emph{overwrite} them in place.
The canonical inner loop looks like this update of a vector $\vy$ by a scaled
vector $\vx$ (the operation BLAS calls \texttt{axpy}, for ``$a$ times $x$ plus
$y$''):

\begin{lstlisting}[style=l4style, language=]
for (int i = 0; i < n; i++)
    y[i] = a * x[i] + y[i];     // overwrite y[i] in place
\end{lstlisting}

\noindent After the loop, the array named \texttt{y} holds different numbers than
before. The \emph{name} \texttt{y} survived; the \emph{values} it points to were
destroyed and replaced. There is exactly one block of memory, and it is mutated.
A whole simulation, in this picture, is a long sequence of such overwrites: the
same arrays, scribbled over again and again, step after step, until the final
state is whatever happens to be left in memory when the loop ends.

This picture is not wrong---it is, more or less, what the silicon does, and it is
why classical simulators are fast. But it has a quiet cost: the meaning of the
program is tangled up with \emph{when} each overwrite happens. To know what
\texttt{y} holds, you must know exactly which loop iterations have run and in what
order. The data and its history are the same thing.

\subsection{The picture we will use: pure functions over immutable values}

Here is the other model. A tensor is a \emph{value}, like the number $7$. You
never change a value; you only compute new ones from it. The expression $3 + 4$
does not ``overwrite'' the $3$---it produces a fresh $7$ and leaves $3$ untouched.
We treat tensors the same way. The \texttt{axpy} update becomes a \emph{pure
function}: it takes $a$, $\vx$, and $\vy$, and \emph{returns a brand-new tensor},
leaving its inputs exactly as they were.

\begin{lstlisting}[style=l4style]
axpy :: Scalar -> Tensor[N] -> Tensor[N] -> Tensor[N]
axpy a x y = a * x + y      -- returns a FRESH tensor; x and y unchanged
\end{lstlisting}

\noindent The arrow signature reads ``\texttt{axpy} takes a scalar and two
length-$N$ tensors and gives back a length-$N$ tensor.'' Nothing is overwritten.
If we want to advance a state $\vs$ repeatedly, we do not loop over memory; we
\emph{thread a sequence of distinct values} through a step function:
\[
  \vs_0 \;\xrightarrow{\;\mathrm{step}\;}\; \vs_1
        \;\xrightarrow{\;\mathrm{step}\;}\; \vs_2
        \;\xrightarrow{\;\mathrm{step}\;}\; \cdots
\]
Each $\vs_{k+1} = \mathrm{step}(\vs_k)$ is a fresh tensor produced from the one
before it. The old states still ``exist'' as values---we simply stop referring to
them, and the machine is free to reclaim their memory. \Cref{fig:mutvalue}
draws the two models side by side.

\begin{figure}[t]
\centering
\begin{tikzpicture}[node distance=9mm]
  % ============ LEFT: in-place mutation ============
  \begin{scope}
    \node[font=\bfseries\small, simRust] at (0,2.55) {In-place mutation};
    % a single memory cell, mutated across steps
    \node[draw=simRust, fill=simBgRust, thick, rounded corners=2pt,
          minimum width=20mm, minimum height=9mm] (cell) at (0,1.3)
      {\ttfamily y};
    \node[font=\scriptsize\itshape, simGray, right=1mm of cell]
      {one block of memory};
    % self-loop "overwrite" arrows
    \draw[backflow] (cell.north west) to[out=140,in=40,looseness=5]
      node[above,font=\scriptsize]{\texttt{y := a*x + y}} (cell.north east);
    \node[font=\scriptsize, simRust] at (0,-0.05) {step 1};
    \node[font=\scriptsize, simRust] at (0,-0.55) {step 2};
    \node[font=\scriptsize, simRust] at (0,-1.05) {step 3};
    \draw[backflow] (-1.0,-0.05) -- (cell.south west);
    \node[align=center, font=\scriptsize\itshape, simGray, text width=34mm]
      at (0,-1.85) {the same cell, scribbled over again and again};
  \end{scope}
  % ============ RIGHT: value threading ============
  \begin{scope}[shift={(6.6,0)}]
    \node[font=\bfseries\small, simBlue] at (0,2.55) {Value threading};
    \node[draw=simBlue, fill=simBgBlue, thick, rounded corners=2pt,
          minimum width=13mm, minimum height=8mm] (s0) at (0,1.4) {$\vs_0$};
    \node[draw=simBlue, fill=simBgBlue, thick, rounded corners=2pt,
          minimum width=13mm, minimum height=8mm] (s1) at (0,0.0) {$\vs_1$};
    \node[draw=simBlue, fill=simBgBlue, thick, rounded corners=2pt,
          minimum width=13mm, minimum height=8mm] (s2) at (0,-1.4) {$\vs_2$};
    \draw[flow] (s0) -- node[right,font=\scriptsize]{\texttt{step}} (s1);
    \draw[flow] (s1) -- node[right,font=\scriptsize]{\texttt{step}} (s2);
    \node[align=center, font=\scriptsize\itshape, simGray, text width=36mm]
      at (0,-2.25) {distinct immutable tensors, each produced from the last};
  \end{scope}
\end{tikzpicture}
\caption[Array mutation versus value threading]{The same computation under two
memory models. \emph{Left:} one block of memory named \texttt{y} is overwritten
in place, step after step; to know what it holds you must know how many steps have
run. \emph{Right:} each step is a pure function producing a \emph{fresh}
immutable tensor $\vs_{k+1}$ from $\vs_k$; the states are distinct values and the
chain of arrows \emph{is} the computation's history.}
\label{fig:mutvalue}
\end{figure}

\begin{pitfall}
The value-threading picture does \emph{not} mean we keep every intermediate
tensor in memory forever, nor that we copy a giant array on every step. It is a
statement about \emph{meaning}, not about machine resources. An implementation is
free to reuse the storage of $\vs_k$ for $\vs_{k+1}$ once nothing else refers to
$\vs_k$---exactly the in-place overwrite, recovered as an optimization. We are
choosing how to \emph{describe} the computation, not forbidding the compiler from
being clever. \Cref{ch:monad} makes this reuse precise.
\end{pitfall}

\subsection{Same algorithm, two memory models}

It is worth seeing, plainly, that these two pictures compute \emph{the same
thing}. Whether we overwrite \texttt{y} in place or thread a chain of fresh
tensors $\vs_0 \to \vs_1 \to \vs_2$, the arithmetic performed is identical and so
is the final answer. \Cref{fig:twomodels} makes this concrete: one dependency
graph, the bare sequence of arithmetic operations, annotated twice---once as
overwrites of a single cell, once as a chain of values.

\begin{figure}[t]
\centering
\begin{tikzpicture}[node distance=11mm]
  % the shared op chain
  \node[opbox, minimum width=14mm] (a) {step};
  \node[opbox, minimum width=14mm, right=14mm of a] (b) {step};
  \node[opbox, minimum width=14mm, right=14mm of b] (c) {step};
  \node[data, left=10mm of a] (in) {field};
  \node[data, right=10mm of c] (out) {result};
  \draw[flow] (in) -- (a);
  \draw[flow] (a) -- (b);
  \draw[flow] (b) -- (c);
  \draw[flow] (c) -- (out);
  % mutation annotation (above, rust)
  \node[font=\scriptsize, simRust] at ($(in)+(0,0.62)$) {\texttt{y}};
  \node[font=\scriptsize, simRust] at ($(a)!0.5!(b)+(0,0.62)$) {\texttt{y}};
  \node[font=\scriptsize, simRust] at ($(b)!0.5!(c)+(0,0.62)$) {\texttt{y}};
  \node[font=\scriptsize, simRust] at ($(c)!0.5!(out)+(0,0.62)$) {\texttt{y}};
  \node[font=\scriptsize\itshape, simRust, anchor=west] at (-3.0,0.62)
    {as mutation:};
  % value annotation (below, blue)
  \node[font=\scriptsize, simBlue] at ($(in)+(0,-0.62)$) {$\vs_0$};
  \node[font=\scriptsize, simBlue] at ($(a)!0.5!(b)+(0,-0.62)$) {$\vs_1$};
  \node[font=\scriptsize, simBlue] at ($(b)!0.5!(c)+(0,-0.62)$) {$\vs_2$};
  \node[font=\scriptsize, simBlue] at ($(c)!0.5!(out)+(0,-0.62)$) {$\vs_3$};
  \node[font=\scriptsize\itshape, simBlue, anchor=west] at (-3.0,-0.62)
    {as values:};
\end{tikzpicture}
\caption[One algorithm, two annotations]{The same dependency graph, annotated
twice. The boxes are the actual arithmetic; the chain of arrows is fixed. Read it
along the top (rust) and every arrow carries the \emph{same name} \texttt{y}---a
single cell, overwritten three times. Read it along the bottom (blue) and every
arrow carries a \emph{distinct value} $\vs_k$---a chain of immutable tensors. The
two readings are the same computation; they differ only in what the names mean.}
\label{fig:twomodels}
\end{figure}

This is the crucial observation, and it is liberating: \emph{choosing the
value-threading reading costs us nothing in the final answer.} It is a change of
\emph{description}, not of algorithm. So if it costs nothing, why prefer it? The
next section is the answer.

\section{Why the pure picture pays}

Three payoffs, in increasing order of how much they shape the rest of the book.

\paragraph{The data flow becomes the dependency graph.}
In the mutation picture, the arrows in \cref{fig:twomodels} are implicit---you
recover them by tracing which loop reads which array after which write. In the
value picture, the arrows are \emph{explicit and literal}: $\vs_1$ is the input to
the step that makes $\vs_2$, full stop. A pure function's output depends on
exactly its inputs and nothing else---no hidden global it quietly read, no array
some other loop overwrote behind its back. So the program text \emph{is} the
mathematical dependency graph. This is not a metaphor; it is what lets us, later,
read off an algorithm's structure---what it iterates, what it folds, what it can
do in parallel---directly from its written form (\cref{ch:fold}).

\paragraph{Pure pieces compose without surprises.}
Because a pure step depends only on its inputs, two steps that do not share an
input cannot interfere, and a step behaves the same no matter what ran before it.
This makes the pieces \emph{composable}: we can snap them together, reason about
each in isolation, substitute one implementation for another, and trust that the
whole behaves as the sum of its parts. Half of this book is the discovery that
six very different-looking simulators are built from one small set of such pieces,
combined in slightly different ways. That discovery is only possible because the
pieces compose cleanly---which is only true because they are pure.

\paragraph{The pure picture is the one the accelerator wants.}
Our eventual target is not a single CPU walking an array one entry at a time; it
is a GPU (or another massively parallel backend) applying an operation to
\emph{all} entries of a tensor at once. Such hardware is happiest when told ``here
is an immutable input tensor; produce an output tensor''---a pure, whole-tensor
transformation with no aliasing and no ordering hazards to police. The
value-threading description is, almost word for word, the form these backends
consume. We call this the \emph{implementation view}: the same spec, read as a
library of pure tensor functions ready to lower onto an accelerator. It is the
destination of the whole book, assembled explicitly in \cref{part:collection}.

\begin{keyidea}
Mutation and value-threading compute the same answer, so the choice is free---and
we spend that freedom on the pure picture because it makes the \emph{data flow
explicit} (the program text is the dependency graph), makes the \emph{pieces
compose} (pure functions don't interfere), and matches the \emph{accelerator
backend} (whole-tensor, no aliasing). Clarity, composability, and the GPU: three
reasons, one decision.
\end{keyidea}

\section{A physical analogy: stream, then collide}

A small physical example makes the ``make a new state from the old one'' rhythm
concrete. The lattice-Boltzmann method simulates a fluid by tracking, at each cell
of a grid, how much fluid is moving in each of a few discrete directions---a
little tensor of \emph{populations} per cell. One time step has two sub-steps. In
\emph{streaming}, every population slides to the neighboring cell in its
direction. In \emph{collision}, the populations within each cell relax toward a
local equilibrium. Stream, then collide; that is one step.

In a classical code this is a nest of loops overwriting a population array in
place. In our picture it is a pure function from the old state to a fresh one:

\begin{lstlisting}[style=l4style]
step :: FluidState -> FluidState
step s =
  let streamed = stream  s.populations in   -- slide to neighbors
  let collided = collide streamed       in   -- relax toward equilibrium
  { ...s, populations: collided, t: s.t + 1 }
\end{lstlisting}

\noindent The step reads the current state \texttt{s}, computes two intermediate
tensors (\texttt{streamed}, then \texttt{collided}), and returns a \emph{new}
record with the advanced populations and an incremented step count. Nothing is
overwritten. A whole simulation is just this step, applied again and again to its
own output---exactly the $\vs_0 \to \vs_1 \to \vs_2$ chain of
\cref{fig:mutvalue}. Our solvers will have the same shape: a step that turns one
field-state into the next, threaded through an iteration. We have borrowed the
fluid's rhythm; we will reuse it for time integration in \cref{ch:transient} and,
combinator by combinator, throughout \cref{part:framework}.

\begin{physicsbox}
``Stream then collide'' is the lattice-Boltzmann heartbeat, but the
\emph{shape}---read the current state, compute a fresh next state, advance a
counter---is universal. A time-domain Maxwell step does it. One iteration of a
linear solver does it. One round of mesh refinement does it. Every sequential
computation in this book is a step function threaded through an iteration; the
physics only changes what happens \emph{inside} the step.
\end{physicsbox}

\section{Three questions the calculus asks}

We have made one decision---describe simulations as pure transformations of
immutable tensors. To turn that decision into a working language, we will, in
\cref{part:framework}, build a small \emph{calculus}: a tiny, precise notation for
exactly these transformations. We close this chapter with a first look at the
three questions that calculus is organized to answer. You have already met all
three informally in this chapter; naming them now plants the signposts for
Part~III.

\begin{enumerate}[label=\textbf{Q\arabic*.}, leftmargin=*]
  \item \textbf{What operations happen?} The verbs---the primitive
  transformations and the named combinations of them. \texttt{axpy} is a verb;
  \texttt{stream} and \texttt{collide} are verbs; ``apply this operator,''
  ``take this inner product,'' ``assemble this matrix'' are verbs. Pinning down
  the vocabulary of operations, each with its shape signature like
  \lstinline[style=l4style]|axpy :: Scalar -> Tensor[N] -> Tensor[N] -> Tensor[N]|,
  is the work of \cref{ch:tensors} and \cref{ch:operators}.

  \item \textbf{Who owns the state?} Not every value plays the same role. The
  evolving field-state is threaded from step to step; a solver's matrix entries
  are fixed once and only read; a step's scratch intermediates live and die
  inside one step. Keeping these straight---which values evolve, which are
  read-only parameters, which are ephemeral---is what lets us thread state purely
  without losing track of it. \Cref{ch:monad} gives this its own discipline,
  the \emph{state monad}, and shows that the entire ``mutation'' of a solve is one
  carefully placed write.

  \item \textbf{How is state evolution coordinated?} A list of verbs is not yet an
  algorithm; something must say \emph{do this, then this, repeat until converged,
  collect the results}. That coordination---sequencing, iteration, convergence
  tests, mapping over a family---is itself made of a few reusable shapes we call
  \emph{combinators}. The most important is the \emph{fold}: apply a step to its
  own output, over and over, threading a state through. \Cref{ch:fold} is built
  around it.
\end{enumerate}

\begin{notationbox}
These three questions---\emph{what operations}, \emph{who owns state}, \emph{how
coordinated}---are exactly the three axes the formal calculus separates, and they
are the spine of \cref{part:framework}: one question per cluster of chapters
(verbs and shapes; the state monad; the fold and its relatives). Whenever a later
chapter feels intricate, return here and ask which of the three it is answering.
\end{notationbox}

We now have the lens. \Cref{ch:situating} uses it immediately: it lays the six
analyses of \cref{ch:targets} onto this framework, showing that the ``solve''
stage of every one is one of just four coordination shapes---four ways of
threading the step. With that, Part~I is complete, and we can begin building the
field theory and the calculus the rest of the book rests on.

\summarysection
A discretized field is a \emph{tensor}: a finite, named-axis array of numbers, the
degrees of freedom collected into a rectangular shape like
\lstinline[style=l4style]|Tensor[N]| or \lstinline[style=l4style]|Tensor[H, W, 3]|.
A simulation is a sequence of transformations of such tensors. There are two ways
to describe ``transformation.'' The classical one---\emph{in-place mutation}---
overwrites a single block of memory step after step, tangling a value with the
history of writes that produced it. The one we adopt---\emph{pure functions over
immutable values}---threads a chain of distinct tensors $\vs_0 \to \vs_1 \to
\cdots$, each freshly produced from the last. The two compute the same answer, so
the choice is free; we spend it on the pure picture because it makes the data flow
into an explicit dependency graph, makes the pieces compose without interference,
and matches the whole-tensor form a GPU/accelerator backend wants (the
implementation view). The lattice-Boltzmann ``stream then collide'' step is the
prototype of the step functions our solvers thread. Finally, the calculus we build
in Part~III is organized around three questions: \emph{what operations happen},
\emph{who owns the state}, and \emph{how state evolution is coordinated}.

\furtherreading
The shift from mutable arrays to pure functions over immutable values is the
central idea of functional programming; the classic, readable entry points are
Wadler's exposition of how monads tame state and effects in a pure
language~\citep{wadler1995monads}, and the standard reference for the language
that crystallized these ideas~\citep{peytonjones2003}. No functional-programming
background is assumed---\cref{ch:monad} builds the state monad we will need from
scratch. The tensor-and-shape model previewed here is developed fully in
\cref{ch:tensors}; the fold, our central combinator, in \cref{ch:fold}. Readers
who want the physical analogy in depth will find the lattice-Boltzmann method in
any computational fluid dynamics text; we use only its rhythm, not its details.

\exercisessection
\begin{enumerate}[nosep]
  \item A scalar potential on a mesh with $N=10^6$ degrees of freedom, and an
  $\vE$-field on the same mesh, are both stored as
  \lstinline[style=l4style]|Tensor[N]|. What does the single axis $N$ count in
  each case, and why is it the \emph{same} $N$? (Hint: revisit which degrees of
  freedom each field lives on; the full story is \cref{ch:functionspaces}.)
  \item Write the \texttt{axpy} update $y \leftarrow a x + y$ once as a C-style
  in-place loop and once as a pure function returning a fresh tensor. For each,
  state precisely what is true of the input \texttt{x} and the input \texttt{y}
  \emph{after} the operation completes.
  \item In \cref{fig:twomodels} the two readings differ only in what the names on
  the arrows mean. Explain why a compiler is nonetheless free to \emph{implement}
  the value-threading reading as in-place mutation. What condition on $\vs_k$ must
  hold for it to reuse that tensor's storage for $\vs_{k+1}$?
  \item Three of this chapter's payoffs were: explicit dependency graph,
  composability, accelerator-friendliness. For each, give a one-sentence reason it
  is a \emph{consequence of purity}---that is, of a step depending on nothing but
  its declared inputs.
  \item The lattice-Boltzmann \texttt{step} returns
  \lstinline[style=l4style]|{ ...s, populations: collided, t: s.t + 1 }|, a fresh
  record. Rewrite ``stream then collide'' as a description that mutates a single
  population array in place, and identify the one place where the pure version's
  meaning would become ambiguous if two such steps were run on overlapping memory.
\end{enumerate}
